\section{阈值极限：重—轻 Sudakov 核与类空喷注函数}\label{sec:yzliu-threshold}

准 PDF 匹配核在 $\xi\to1$ 处含 $\ln^k|1-\xi|/(1-\xi)$ 型的大对数。Ji、Liu 与 Su\cite{Ji:2023pba}论证了它们与 $x\to1$ 的 DIS 一样可以因子化（\cref{fig:yzliu-threshold}）：硬核正是准 TMD 的重—轻核 $H$，半硬标度 $|1-\xi|p^z$ 处的交换由一个类空喷注函数重求和。

%%FIG fig:yzliu-threshold

\begin{definition}[重—轻 Sudakov 因子与类空喷注函数]\label{def:yzliu-hl-jet}
(a) $C^\pm_{HL}$ 为单个顶点的重—轻 Sudakov 因子（在壳轻夸克变为静态类空连线）；文献\cite{Ji:2023pba}的 “$H_{HL}$” 指一对顶点；$H\equiv|H_{HL}|=C^+_{HL}C^-_{HL}$（\cref{prop:yzliu-hard-kernel}）。(b) 类空喷注函数 % VF-LAM-14, VF-YZL-9
\begin{equation}\label{eq:yzliu-jet-def}
\tilde J(\mu z)=\langle\Omega|\mathcal T\,W_{n,+}(zn_z)\,W_z^\dagger(zn_z)\,W_z(0)\,W_{n,-}(0)|\Omega\rangle ,
\end{equation}
$W_{n,-}(0)$ 为自 $-\infty n$ 到 $0$ 的类光连线，$W_{n,+}(zn_z)$ 为自 $zn_z$ 到 $+\infty n$ 的类光连线，$W_z^\dagger(zn_z)W_z(0)$ 为 $0$ 与 $zn_z$ 之间的类空连线（\cref{def:prelim-wilson}）：类光静态色源变为类空静态色源再变回的“跃迁形状因子”。它在阈下、为实、只依赖 $|z|$，共轭于半硬标度 $|1-\xi|p^z$；类时版本 $\tilde J(\mu t)$ 由 $n_z\to n_t$ 得到。(c) $J_f(z\mu)\equiv\tilde J(z^2\mu^2)\,\ee^{\ii\,\mathrm{sgn}(z)A(L_\zeta)}$，$\mathrm{sgn}(z)A=\mathrm{Arg}\,H_{HL}$。约定：我们的算符 $\bar\psi(z)\gamma^0W[z,0]\psi(0)$（\cref{def:lattice-quasi}）与文献\cite{Ji:2023pba}相同，$\mathrm{sgn}(z)$ 含义一致；该文的 $\lambda=zp^z=-\lambda_{\text{本书}}$，其 $L_z$、$l_z$ 即本书的 $L_\zeta$、$L_z=\ln(z^2\mu^2\ee^{2\gamma_E}/4)$。
\uses{def:prelim-wilson,def:lattice-quasi,prop:yzliu-hard-kernel}
\end{definition}

\begin{theorem}[类空与类时重夸克喷注函数的解析延拓 \tagA]\label{thm:yzliu-jet-continuation}
在维数正规化微扰论各阶，对 $t>0$，$\tilde J(t=-\ii|z|,D)=\tilde J(z,D)$。因此两圈类空喷注函数可由类时重夸克喷注函数\cite{Jain:2008gb}得到。单圈 % VF-YZL-9
\begin{equation}\label{eq:yzliu-jet-1loop}
\tilde J=1+a\Big(\frac12L_z^2+L_z+\frac{\pi^2}{12}+2\Big),\qquad \frac{\dd\ln\tilde J}{\dd\ln\mu}=\Gamma_{\mathrm{cusp}}L_z-\tilde\gamma_J,\qquad \tilde\gamma^{(1)}_J=-\frac{\alpha_sC_F}{\pi},
\end{equation}
两圈 $\tilde\gamma^{(2)}_J=C_FC_A(-\frac{1396}{27}+\frac{23\pi^2}9+20\zeta_3)+C_Fn_f(\frac{232}{27}-\frac{2\pi^2}9)$（以 $(\frac{\alpha_s}{4\pi})^2$ 为单位）。（本定理只涉及真空中的 Wilson 线关联函数。）
\uses{def:yzliu-hl-jet,asm:yzliu-wightman-pt}
\end{theorem}

\begin{proofsketch}
证明一\cite[附录~C.1]{Ji:2023pba}：在坐标空间用胶子 $m$ 点关联函数的 Schwinger 表示，缩并保持 $\int D\alpha\,\mathcal P(\alpha)\exp[-\frac\ii2\sum Q_{ij}(x_i-x_j)^2]$、$Q\ge0$ 的形式；对剩余的标度参数 $s$，所有留数极点位于第二象限，故可转动 $s\to-\ii s$，再令 $t\to-\ii z$；自旋结构逐项对应\cite[式~(C.19)--(C.21)]{Ji:2023pba}。证明二\cite[附录~C.2]{Ji:2023pba}：前向管中的 Wightman 函数加复 Lorentz 协变性（\cref{asm:yzliu-wightman-pt}）。单圈：裸表达式在 $\overline{\mathrm{MS}}$ 中的极点为 $a[\frac1{\varepsilon^2}+\frac{L_z+1}{\varepsilon}]$，有限部分即\cref{eq:yzliu-jet-1loop}；$\frac{\dd L_z}{\dd\ln\mu}=2$ 给出 $\frac{\dd\tilde J^{(1)}}{\dd\ln\mu}=2a(L_z+1)$，即 $\Gamma^{(1)}_{\mathrm{cusp}}=2a$、$\tilde\gamma^{(1)}_J=-2a$。两圈：$\tilde\gamma_J=2\gamma_{HL}-2\gamma_s$，由\cite[式~(A.4)--(A.5)]{Ji:2023pba}得上式，与类时重夸克喷注函数的反常维度相同。
\end{proofsketch}

\paragraph{地位与缺口。}全阶论证 [A]，与\cref{thm:yzliu-time-ordering}同类，证明一给出完整的组合概要，证明二的瓶颈是\cref{asm:yzliu-wightman-pt}。$\tilde\gamma_J$ 等于类时重夸克喷注函数的反常维度；$\tilde J$ 为实的论证：在时序微扰论中类光连线的能量分母同号，唯一的虚部来源是 $t=0$ 的静态类空连线，而 $z>0$ 时回路只含类空与类光间隔、处在阈下。

\begin{theorem}[准 PDF 匹配核的阈值因子化（Ji--Liu--Su） \tagA]\label{thm:yzliu-threshold}
对\cref{thm:lattice-lamet}的非单态 $\overline{\mathrm{MS}}$ 核（$\Gamma=\gamma^0$，$\xi=x/y$），在 $1-\xi$ 的领头幂次（所有 $\propto(1-\xi)^{-1}\ln^k|1-\xi|$ 与 $\delta(1-\xi)$ 的项），微扰论各阶
\begin{equation}\label{eq:yzliu-threshold}
C\Big(\xi,\frac{p^z}{\mu}\Big)\Big|_{\xi\to1}=H\Big(\frac{4p_z^2}{\mu^2}\Big)\,p^z\,J_f\Big(\frac{(1-\xi)p^z}{\mu},\frac{4p_z^2}{\mu^2}\Big)\big[1+\order{1-\xi}\big],\qquad p^z=|y|P^z ,
\end{equation}
$H=|H_{HL}|$ 与准 TMD、准光前波函数因子化中的核相同（普适性），$J_f$ 如\cref{def:yzliu-hl-jet}。单圈：$H^{(1)}=a(-2+\frac{\pi^2}{12}-\frac{L_\zeta^2}2+L_\zeta)$，$\tilde J^{(1)}$ 如\cref{eq:yzliu-jet-1loop}，$\Imag J^{(1)}_f=a\pi\,\mathrm{sgn}(z)(L_\zeta-1)$；在动量空间 % VF-LAM-14
\begin{equation}\label{eq:yzliu-Jf-mom}
p^zJ_f=\delta(1-\xi)\Big[1+a\Big(2+\frac{\pi^2}4+\frac{L_\zeta^2}2-L_\zeta\Big)\Big]+a\Big[\mathcal P\frac{\ln(1-\xi)^2}{|1-\xi|}+(L_\zeta-1)\Big(\mathcal P\frac1{1-\xi}+\mathcal P\frac1{|1-\xi|}\Big)\Big].
\end{equation}
RGE：$\frac{\dd\ln H}{\dd\ln\mu}=\Gamma_{\mathrm{cusp}}L_\zeta+\tilde\gamma_H$，$\frac{\dd\ln\tilde J}{\dd\ln\mu}=\Gamma_{\mathrm{cusp}}L_z-\tilde\gamma_J$，$\frac{\dd\ln H_{HL}}{\dd\ln\mu}=\Gamma_{\mathrm{cusp}}L_\zeta+\tilde\gamma_H-\ii\pi\,\mathrm{sgn}(z)\Gamma_{\mathrm{cusp}}$；单圈时总反常维度 $\Gamma_{\mathrm{cusp}}[\ln(\lambda^2\ee^{2\gamma_E})-\ii\pi\,\mathrm{sgn}z]$ 与 $\mu$ 无关。（本定理只涉及部分子态。）
\uses{def:yzliu-hl-jet,thm:yzliu-jet-continuation,prop:yzliu-hard-kernel,prop:lamet-oneloop,thm:lattice-lamet}
\end{theorem}

\begin{proofsketch}
（1）领头区域\cite[第~1--2~节、图~1]{Ji:2023pba}：入射与出射共线夸克在两个夸克—连线连接点处有硬顶点（标度 $\zeta_z=4p_z^2$），半硬胶子（$|1-\xi|p^z$）在共线夸克与静态连线之间交换——相当于 $x\to1$ 的 DIS\cite{Sterman1987,BecherNeubertPecjak2007}，类空连线取代了反共线夸克；两次 $k_z=\order{p_z}$ 相反的硬实交换受相空间压低。（2）$\tilde J$ 是真空 Wilson 回路关联函数，经\cref{thm:yzliu-jet-continuation}与类时重夸克喷注函数相联。（3）RG 相容性再现 DGLAP 的阈值极限（见下面的反常维度代数）。（4）定阶：单圈精确成立（与\cref{prop:lamet-oneloop}比较，见下例）；把 $H^{(1)}$、$\tilde J^{(1)}$ 与相位代入，$\frac{\dd}{\dd\ln\mu}[H^{(1)}+\tilde J^{(1)}+\ii\pi\,\mathrm{sgn}(z)a(L_\zeta-1)]=2a[L_\zeta+L_z-\ii\pi\,\mathrm{sgn}z]$，而 $L_\zeta+L_z=\ln(\lambda^2\ee^{2\gamma_E})$。NNLO 检验见下文。
\end{proofsketch}

\paragraph{地位与缺口。}全阶论证 [A]：领头区域论证加时空图像与 RG 相容性（SCET 型），不是完整的 Collins 式幂次计数证明；相对于 DIS 阈值因子化（\cref{prop:thr-dis}），新的要素是相位。单圈的坐标空间核在 $|\lambda|$ 直到 $10^5$ 处与\cref{eq:yzliu-threshold}一致，包括 $\lambda<0$ 时相位的符号。归属：文献\cite{Ji:2023pba}推导了动量空间核的阈值因子化 $C\to H\cdot J$ 及其 RG 重求和；更早，文献\cite{Gao:2021hxl}解释了阈值对数的来源，并在 NLL 重求和了赝 PDF 系数中的阈值对数。
% Verification: claim N06 (one-loop check to |lambda| = 1e5, verification, no marker); script s08 (W26).

\begin{example}[阈值因子化的单圈检验]\label{ex:yzliu-threshold-oneloop}
$0<\xi<1$：$\overline{\mathrm{MS}}$ 准 PDF 核的奇异部分为 $\frac{2a}{1-\xi}(L_\zeta+\ln(1-\xi)-1)$（\cref{prop:lamet-oneloop}(e)）；由 $\mathcal P\frac1{|1-\xi|}+\mathcal P\frac1{1-\xi}=\frac2{1-\xi}$，它等于\cref{eq:yzliu-Jf-mom}的 $J^{(1)}_f$。$\xi>1$：两个主值相消，核给出 $\frac{2a\ln(\xi-1)}{\xi-1}=J^{(1)}_f$。$\delta(1-\xi)$ 项的系数给出 $H^{(1)}$，与\cref{prop:yzliu-hard-kernel}(b) 的准 TMD 核相同。坐标空间中，记单圈的坐标空间核为 $h(\lambda)$（$\lambda=p\cdot z$），则 $\ee^{-\ii\lambda}h(\lambda)-1-[H^{(1)}+\tilde J^{(1)}+\ii\pi a\,\mathrm{sgn}(z)(L_\zeta-1)]$ 在 $|\lambda|\to\infty$ 时按 $\ln^2|\lambda|/|\lambda|$ 趋于零。
\uses{thm:yzliu-threshold,prop:lamet-oneloop}
\end{example}

\begin{proposition}[反常维度代数 \tagA]\label{prop:yzliu-ad-algebra}
在基 $\{\Gamma_{\mathrm{cusp}},\gamma_F,\gamma_V,\gamma_s,\gamma_{HL}\}$ 下，% VF-YZL-10
\begin{equation}\label{eq:yzliu-ad-algebra}
\begin{gathered}
\gamma_\phi=-\tfrac12\gamma_V-\tfrac12\gamma_s,\qquad \gamma^{\mathrm{thr}}_J=\tfrac12\gamma_V-\tfrac12\gamma_s,\qquad \tilde\gamma_J=2\gamma_{HL}-2\gamma_s,\\
\gamma_{HH}=2\gamma_{HL}-\gamma_s,\qquad \tilde\gamma_H=\gamma_V+2\gamma_F+2\gamma_{HL}-\gamma_s ;
\end{gathered}
\end{equation}
阈值处的 DGLAP 相容性 $2\gamma_\phi+\tilde\gamma_H-\tilde\gamma_J=2\gamma_F$（端点形式 $P_{qq}\to\frac{2\Gamma_{\mathrm{cusp}}}{(1-z)_+}+2\gamma_\phi\delta(1-z)$）。（两圈内容，\tagR，$R_{\mathrm{PT}}$）（$R_{\mathrm{PT}}$：在重整化微扰论中对红外正规化的部分子态严格。）文献\cite{Ji:2023pba}的两圈输入 (A.1)--(A.5) 再现 $\gamma^{(2)}_{HH}$（独立地见\cref{prop:yzliu-sI-twoloop}）与 $\tilde\gamma^{(2)}_H$（\cref{prop:yzliu-hard-kernel}(c)）；$\gamma_V$ 为标准的夸克形状因子非尖点反常维度。（本命题只涉及部分子态。）
\uses{thm:yzliu-threshold,prop:yzliu-hard-kernel,thm:yzliu-intrinsic-soft,thm:yzliu-jet-continuation}
\end{proposition}

\begin{proof}
对\cref{eq:yzliu-threshold}取 $\frac{\dd}{\dd\ln\mu}$：$H$ 贡献 $\tilde\gamma_H$，$J_f$ 在动量空间满足 $\frac{\dd J_f}{\dd\ln\mu}=-(\tilde\gamma_J+\Gamma_{\mathrm{cusp}}\ln\frac{4p_z^2}{\mu^2})J_f-2\Gamma_{\mathrm{cusp}}\mathcal P\!\int_{\xi'>\xi}\frac{J_f(\xi')}{\xi'-\xi}\dd\xi'$，于是 $C$ 的演化含 $(\tilde\gamma_H-\tilde\gamma_J)C$ 与一个积分项\cite[附录~A]{Ji:2023pba}。积分项再现 DGLAP 核的 $\frac{2\Gamma_{\mathrm{cusp}}}{(1-z)_+}$；$\delta(1-z)$ 项相符要求 $2\gamma_\phi+\tilde\gamma_H-\tilde\gamma_J=2\gamma_F$（$\gamma_F$ 来自准 PDF 的乘法重整化，\cref{prop:yzliu-hard-kernel}(a)）。DIS 阈值因子化给出 $\gamma_V+(\gamma_\phi+\gamma_s)+\gamma_\phi=0$，其中 $\gamma_\phi+\gamma_s=-\gamma^{\mathrm{thr}}_J$，故 $2\gamma_\phi=-\gamma_V-\gamma_s$。代入 $\tilde\gamma_J=2\gamma_{HL}-2\gamma_s$ 即得 $\tilde\gamma_H=\gamma_V+2\gamma_F+2\gamma_{HL}-\gamma_s$；反之，把\cref{eq:yzliu-ad-algebra}代入，$2\gamma_\phi+\tilde\gamma_H-\tilde\gamma_J=(-\gamma_V-\gamma_s)+(\gamma_V+2\gamma_F+2\gamma_{HL}-\gamma_s)-(2\gamma_{HL}-2\gamma_s)=2\gamma_F$ 恒成立。$\gamma_{HH}=2\gamma_{HL}-\gamma_s$ 由角依赖尖点反常维度的大角渐近 $\Gamma_{\mathrm{cusp}}(\alpha_s,Y)\to Y\Gamma_{\mathrm{cusp}}+\gamma_{HH}$ 与\cref{thm:yzliu-intrinsic-soft}(d) 得到，$\tilde\gamma_J$ 的结构见\cref{thm:yzliu-jet-continuation}。两圈：把 (A.1)--(A.5) 代入，得 $\gamma^{(2)}_{HH}$ 与上面所引的两圈值一致，得 $\tilde\gamma^{(2)}_H$ 与\cref{prop:yzliu-hard-kernel}(c) 的\cref{eq:yzliu-gH2}一致（即 (A.11) 与 (A.13) 相符）；单圈 $\gamma^{(1)}_\phi=\frac{3\alpha_sC_F}{4\pi}$。
\end{proof}

\paragraph{地位与缺口。}结构关系为全阶论证 [A]（来自 [A] 级因子化公式的相容性）；两圈值为 [R, $R_{\mathrm{PT}}$]。

\begin{proposition}[阈值因子化的 NNLO 检验 \tagP]\label{prop:yzliu-threshold-nnlo}
精确的两圈非单态核\cite{Li:2020xml,Chen:2020ody}（\cref{prop:lamet-twoloop}）的阈值极限再现\cref{thm:yzliu-threshold}的全部 41 个非零奇异项（$\delta(1-\xi)$ 与 $\ln^k(1-\xi)/(1-\xi)$，$k\le3$），其中两个常数由比较定出：$c_a$（两圈相位）与 $c_H$（两圈 $\ln H$，\cref{prop:yzliu-hard-kernel}(c)）；独立计算的 NNLO 准 TMD 系数\cite{delRio:2023pse}与此 $c_H$ 一致。（本命题只涉及部分子态。）
\uses{thm:yzliu-threshold,prop:yzliu-hard-kernel,prop:lamet-twoloop}
\end{proposition}

\paragraph{地位与缺口。}[P]：带两个拟合常数的定阶验证\cite[第~5~节、式~(5.7)、(5.10)--(5.11)]{Ji:2023pba}。

\begin{remark}[阈值重求和与大 $x$ 适用域]\label{rem:yzliu-threshold-domain}
重求和核\cite[式~(6.19)]{Ji:2023pba}（$\mu_i=|1-\xi|p^z$，$\eta=2a_\Gamma(\mu_i,\mu)$，$\hat A=A(\alpha(\zeta_z))+\pi a_\Gamma(\zeta_z,\mu)$，$[\ ]_\star$ 为该文式~(6.18) 的星分布）为 % TR, checked on the rendered page 18 of arXiv:2305.04416v1 (eqs. (6.13)-(6.19); eta, A-hat, mu_i as defined there)
\begin{equation*}
\begin{aligned}
C={}&H(\alpha(\zeta_z))\,\ee^{2S(\zeta_z,\mu)-2S(\mu_i,\mu)-a_H(\zeta_z,\mu)+a_J(\mu_i,\mu)}\,\tilde J(L_z=-2\partial_\eta,\alpha(\mu_i))\\
&\times\Big[\frac{\sin(\mathrm{sgn}(1-\xi)\hat A+\frac{\eta\pi}2)}{|1-\xi|}\Big(\frac{2|1-\xi|p^z}{\mu_i}\Big)^\eta\Big]_\star\frac{\Gamma(1-\eta)\ee^{-\eta\gamma_E}}\pi ;
\end{aligned}
\end{equation*}
微扰喷注函数要求 $\Lambda\ll|1-\xi|p^z\ll p^z$。数值上（据\cite{Ji:2024hit}的摘要与引言；$x$ 为光锥动量分数）：硬部分在 $xP^z$、阈值软函数在 $(1-x)P^z$，$\alpha_s^k\ln^{2k}(1-x)$ 的重求和“至关重要”；$(1-x)P^z\to\Lambda$ 时匹配失效，$xP^z$ 与 $(1-x)P^z$ 均远大于 $\Lambda$ 时可用；线性重整子重求和可与阈值重求和结合。更早的 Mellin 空间 NLL 见\cite{Gao:2021hxl}，GPD 的阈值重求和见\cite{Holligan:2025baj}。本注细化了\cref{rem:lamet-domain}。
\end{remark}
